\chapter{Identity: Control, Mutability, and the Counter Defect}

\section{The premise the implementation lacked}

The founding design document opens its universal-substrate section with a sentence that reads like a statement of fact about the running system. Until 2026-07-17 it was not. It was a statement of intent, and the system did not satisfy it. Here it is, verbatim, from \texttt{DESIGN.md} line 100, added 2026-07-04:

\begin{quote}
\textbf{Identifiers.} Identity is a rigid designator (o.i, Asana gid, DOI, git sha), minted at spawn, immutable while state evolves. Names are witnesses attached to IDs, never identities.
\end{quote}

And here was the schema it described, from the Names instance's \texttt{schema.sql}:

\begin{lstlisting}
CREATE TABLE o (
  i BIGINT GENERATED BY DEFAULT AS IDENTITY PRIMARY KEY,
  j JSONB NOT NULL
);
\end{lstlisting}

\texttt{GENERATED BY DEFAULT AS IDENTITY} is a sequence. A sequence is a counter, and a counter is local to the Postgres instance that holds it. Boot a second copy from a nightly dump, insert an object on each, and both mint the same number. The number designates two different things, and nothing in the data says so.

A rigid designator, in Kripke's sense and in the sense \texttt{DESIGN.md} plainly intends, denotes the same object in every world in which that object exists. The worlds here are copies. \texttt{o.i} denoted different objects in different copies. It was therefore not a rigid designator; it was a rigid designator relative to one machine, which is a different and much weaker property, and it is the property that a single-machine system cannot distinguish from the real one. The founding document asserted as a premise the exact property the implementation lacked, and it asserted it while there was only one machine, where the difference is invisible.

The consequence, stated in its sharpest form, because it is easy to read the defect as a merge inconvenience rather than as an expressive gap: \textbf{``the same object on two copies'' was not expressible.} There was no term in the language of the database that referred to an object across copies. \texttt{o.i = 113403} was a sentence whose truth conditions depended on which container you were talking to, and the sentence carried no marker saying so.

That is the past tense throughout, and it is earned. The premise is now discharged.

\section{The fix, and how it is checked}

The substrate that boots today is the stem cell, \texttt{stem-cell.sql}, and it mints as follows:

\begin{lstlisting}
CREATE TABLE o (
  i UUID PRIMARY KEY DEFAULT uuidv7(),
  j JSONB NOT NULL
);

CREATE TABLE r (
  i  UUID NOT NULL DEFAULT uuidv7(),
  o1 UUID,  o2 UUID,  d JSONB NOT NULL,
  c  REAL NOT NULL DEFAULT 1,
  b  UUID NOT NULL DEFAULT gen_random_uuid(),
  t  TIMESTAMPTZ NOT NULL DEFAULT now()
);
\end{lstlisting}

This is exactly the whitepaper's prescription, split along the taxonomy's own seam. Content gets its fingerprint: \texttt{uuid\_from\_sha256()} in \texttt{differentiate-instrument.sql} turns a sha256 into a UUID, and \texttt{blob()} uses it as the primary key, so two machines that capture the same bytes mint the same object without talking. Everything else gets a random, time-ordered identifier: \texttt{uuidv7()}, implemented in the file because Postgres 17 does not have one, verified again on the live container on 2026-07-17 (\texttt{PostgreSQL 17.10}; the built-in landed in Postgres 18). The implementation takes a \texttt{gen\_random\_uuid()}, overlays the first 48 bits with the millisecond clock, and flips the version nibble from 0100 to 0111 by setting bits 52 and 53.

Time-ordering is not decoration here and the stem cell says why: \texttt{r} is a hypertable on \texttt{t} with time-ordered indexes, so v4 keys would scatter every write across the whole keyspace, and v7 keeps them local. That is an index-locality argument, not an identity argument, and it should be read as one.

\begin{Correspondence}{\texttt{o.i} is now a global name, and the check is two containers that never spoke}{Boot two instances from the same stem cell, insert distinct samples and one identical file, and compare the ids.}
This is the check the old chapter said would settle it, and it has now been run rather than reasoned about. Two throwaway containers, same image, same two SQL files, no network path between them and no shared state beyond the text of the schema. Distinct objects got distinct ids; the same file got the same id on both.
\begin{Check}
-- On each container, from stem-cell.sql + differentiate-instrument.sql:
SELECT instance_id();
--   A: 019f6e52-c691-7d0e-914e-9ff0bfaecea6
--   B: 019f6e52-cd8c-767f-a799-1f7c224cb9d1     distinct, uncoordinated

INSERT INTO o (j) VALUES ('{"kind":"sample","id":"LAO-STO-0417B"}') RETURNING i;
--   A: 019f6e52-f2df-707f-9b75-9c74c2394306
--   B: 019f6e52-fb5d-7bf5-8437-8c60e3ed05ae     distinct: two rows, two names

SELECT blob('35e3bc005f4399f8ac0919177ed2275c54b24cff170b833141f815cdc0823c37');
--   A: 35e3bc00-5f43-89f8-ac09-19177ed2275c
--   B: 35e3bc00-5f43-89f8-ac09-19177ed2275c     identical: one file, one name
\end{Check}
The old text ran the mirror image of this check and reported that both sides returned the same integer for two different objects. Both halves now come out the way the design says they should, and neither container was told the other existed.
\end{Correspondence}

\section{The v8 id is an index, not the evidence}

One detail in \texttt{uuid\_from\_sha256()} is worth stating plainly, because it is the sort of thing that gets read as a hash and is not one. The output is UUID version 8, which RFC 9562 reserves for custom application-defined layouts, and its payload is the leading 128 bits of the sha256 with six of them overwritten: four for the version nibble, two for the variant. So the id is a 122-bit truncation of the digest.

A truncation is not the digest. It is a name that is a function of the content, which is all the merge argument needs, and it is not evidence about the content, which is what a fingerprint is for. The stem cell keeps the two roles apart: the full 64 hex characters stay in \texttt{j->>'id'}, and that is the string you verify bytes against. The UUID is the index.

\texttt{blob()} enforces the one invariant that makes any of this mean something. It rejects anything that is not 64 hex characters, so a caller cannot smuggle in a non-fingerprint and have it treated as one, and it builds the row as \verb@meta || jsonb_build_object(...)@ rather than the reverse, so the computed identity fields win over caller-supplied \texttt{meta} instead of losing to it. The old order let a caller pass \verb|{"id":"NOT-THE-HASH"}| and mint a ref whose recorded id was not the hash of anything. Checked on a throwaway container: the forged \texttt{id} is now discarded and the row records the real digest; \texttt{blob('deadbeef')} raises.

\begin{Analogy}{The v8 id is a fingerprint}{Breaks because six bits are structural, so it is 122 bits of digest, and it is a name rather than a proof.}
For collision resistance against an adversary the difference matters: 122 bits truncated from sha256 is not sha256, and nothing in the design asks it to be. What the id has to do is be a deterministic function of the content, so that two uncoordinated machines agree, and be usable as a Postgres primary key. It does both. If you want to know whether these bytes are those bytes, hash the bytes and compare against \texttt{j->>'id'}. Confusing the index with the evidence is how a content-addressed store quietly becomes a store addressed by something that used to be content.
\end{Analogy}

\section{An instance is identified by its big bang}

The identity question does not stop at rows. If two copies are going to merge, each has to be able to say which one it is, and the obvious implementations are a configuration file holding a name or a column holding an instance id. The stem cell does neither.

Every instance records the moment it was brought to life, once, at boot, as an ordinary object with \texttt{kind:'big-bang'} carrying \texttt{at} and \texttt{t\_universe}. That row already has an identity: its own \texttt{o.i}, a uuidv7, unique across every machine without coordination. So the instance is the event, and \texttt{instance\_id()} is a query:

\begin{lstlisting}
CREATE OR REPLACE FUNCTION instance_id() RETURNS UUID LANGUAGE sql STABLE AS $$
  SELECT i FROM o WHERE j->>'kind' = 'big-bang' ORDER BY i LIMIT 1
$$;
\end{lstlisting}

This is the same move the design makes everywhere else, and it is worth naming as such rather than treating it as a trick. Persons are not stored; they are connected components of \texttt{same} witnesses. Truth is not stored; it is a view. The instance is not stored either; it is computed from the one thing that was witnessed about it, namely that it was born. A config file would have been a second source of truth about identity, mutable, editable, copyable along with the volume, and silently wrong on the copy. The birth event cannot be copied without copying its id, and two instances restored from the same dump would share a big bang, which is the correct answer: they are the same instance, forked, and the design has a word for that.

The live instance, queried 2026-07-17: \texttt{instance\_id()} returns \texttt{019f6e4d-adb0-7c46-9d4b-ce4fbe666174}, born \texttt{2026-07-17T04:20:12.336Z}, holding ten objects and zero relations.

\begin{Correspondence}{The instance is not stored, it is computed from its birth}{Ask an instance who it is and note that the answer is a \texttt{SELECT} over \texttt{o}, not a setting.}
\begin{Check}
SELECT instance_id();   -- 019f6e4d-adb0-7c46-9d4b-ce4fbe666174, the big bang's own o.i
SELECT age();           -- seconds since it, a difference, therefore exact
SELECT count(*) FROM o WHERE j->>'kind' = 'big-bang';   -- 1, and it is inserted once, at boot
\end{Check}
\end{Correspondence}

\section{The taxonomy's broken cell, and the repair}

\texttt{DESIGN.md} line 100 continues into the control $\times$ mutability taxonomy, which Chapter 17 argues is the most successful classification in the corpus. Four cells:

\begin{itemize}
\item controlled + immutable: content-addressed (sha256 is the ID)
\item controlled + mutable: ordinary CDB object (\texttt{o.i} + append-only relations)
\item external + immutable: boundary ref \verb|{kind:'ref', scheme, id}|
\item external + mutable: snapshot-on-observation (dated, hashed)
\end{itemize}

Three of the four cells always named an identity mechanism that worked. \texttt{sha256} is a function of content and is the same everywhere by construction. A DOI is minted by Crossref, which is one authority, so it is global. An Asana gid is minted by Asana, likewise. A git sha is a hash. Every one of those names survives being carried to another machine.

The one that did not was the cell the system is made of. ``controlled + mutable: ordinary CDB object'' held all 178{,}561 objects and all 256{,}619 relations of the Names instance, and it was the only cell whose identity mechanism was a local counter. The taxonomy was right that control is the axis that decides representation. What it did not notice is that control is not binary, is not permanent, and is exactly what the lab transition changes. Jeremy controlled \texttt{o.i} because there was one writer. The moment a student's instrument PC writes to its own copy, ``controlled'' means ``controlled by the group,'' and two machines both control the same integer space and mint into it independently.

So the defect was not a corner of the design. It was the diagonal entry of the design's own 2$\times$2. \texttt{uuidv7()} closes that cell by giving the controlled-mutable case a name that needs no authority and no coordination, which is what the other three cells had all along.

\section{Fourteen months, and nobody was ever forced to look}

The counter was not a recent slip. It was the oldest line of code in the project, and it was visible in the first schema block ever written.

\texttt{main.tex} (2025-05-28), the earliest draft in the Categorical Database Architecture folder, line 106: \texttt{item\_id SERIAL PRIMARY KEY}. Line 117: \texttt{correlation\_id SERIAL PRIMARY KEY}. Every one of the four drafts in that folder inherits it by find-and-replace. \texttt{CDA.tex} (2025-05-29), lines 165 and 176: \texttt{object\_id SERIAL PRIMARY KEY}, \texttt{arrow\_id SERIAL PRIMARY KEY}. \texttt{Category Theory for Physicists} does the same with bare \texttt{INTEGER} columns: \texttt{source\_object INTEGER}, \texttt{target\_object INTEGER}, \texttt{observer\_id INTEGER}, \texttt{morphism\_id INTEGER REFERENCES arrows}. Not one of these documents remarks on it.

And the same drafts list, under Future Work, verbatim: ``Distributed Implementation: Scaling atomic architectures across multiple database instances.'' \texttt{CDA.tex} lists ``Cross-Domain Integration: Extending observer relationships to span multiple research domains and institutional boundaries.'' The papers propose distribution in one section and make it impossible in another, four pages apart, with no cross-reference and no awareness.

The reason this survived every rewrite is the useful part, and it is a claim about method rather than about Postgres. The project's rewrites were driven by mathematics: three tables became two, arrows became reifiable, entities stopped being stored, the composition \texttt{CASE} arm was deleted in favor of computed views. Every one of those changes was forced by asking what the structure required. \textbf{Nothing in the mathematics ever forced anyone to ask what a name meant on a second machine}, because category theory does not care how you spell an identifier. The defect lived in the one layer that the project's own strongest instrument does not illuminate.

That the fix took an afternoon once it was written down is not evidence against any of this. Fourteen months is the measurement, and it is a measurement of the method, not of the difficulty.

\section{The reverse finding: the branch column got it right}

There is a genuinely awkward wrinkle here, and it should be reported rather than smoothed. The 2025 branching paper (\texttt{Spacetime Database Topology/main.tex}, whose folder name is a lie and whose real subject is quantum mechanics from classical branching) was already using UUIDs for branch identity. Its Methods say, verbatim: ``Branches indexed by UUID, confidence values stored as REAL.''

The \texttt{b} column has been \texttt{UUID NOT NULL DEFAULT gen\_random\_uuid()} throughout and was always fine: globally unique with no coordination, exactly as it should be. It was \texttt{i} that was a counter. So on the face of it the project used UUIDs where identity was hard and regressed to integers where identity was easy, which is backwards. Note that \texttt{b} is still \texttt{gen\_random\_uuid()} in the stem cell and not \texttt{uuidv7()}, which is right: a branch id is a label with no ordering claim to make.

\begin{Analogy}{The branch column shows the project once knew better}{Breaks because \texttt{b} and \texttt{i} solve different problems, and the UUID choice for \texttt{b} was not a considered position on distribution.}
A branch is minted by a writer who has nothing to coordinate with and no reason to want a small integer; \texttt{gen\_random\_uuid()} is simply the default way to get an opaque token in Postgres, and it is what you reach for when you do not want the value to mean anything. \texttt{i} is a primary key, and the default way to get a primary key in Postgres is a sequence. Both columns took the path of least resistance, and the paths happened to diverge. Reading foresight into \texttt{b} is reading a coin flip as a decision. What survives is the weaker and still useful fact: the mechanism was available, in this project's own files, and was not applied where it mattered.
\end{Analogy}

\section{The second defect: \texttt{r.i} is not unique and still cannot be}

This one was never in the whitepaper or in \texttt{DESIGN.md}, and it is the part of this chapter the big bang did not touch. \texttt{uuidv7()} fixed the names. It did not fix the enforcement, because the enforcement mechanism is unavailable rather than unused.

\texttt{o.i} is a \texttt{PRIMARY KEY}. \texttt{r.i} is not. Query \texttt{pg\_constraint} against the live instance, born four minutes earlier and running the current stem cell, and the answer is the same as it was under the counters:

\begin{lstlisting}
select conname, contype from pg_constraint where conrelid = 'r'::regclass;
-- (0 rows)

select conname, contype from pg_constraint where conrelid = 'o'::regclass;
-- o_pkey | p
\end{lstlisting}

\texttt{r} has no constraints at all. Not a primary key, not a unique index, not a foreign key. \texttt{r.i} carries \texttt{NOT NULL} and a non-unique btree index \texttt{r\_i\_idx}, which is an access path and enforces nothing.

The cause is structural. TimescaleDB requires that any unique constraint on a hypertable include the partitioning column. The partitioning column is \texttt{t}. So the only unique constraints available on \texttt{r} are on tuples containing \texttt{t}, and \texttt{(i, t)} being unique does not make \texttt{i} unique. \texttt{r.i} cannot be a primary key while \texttt{r} is a hypertable on \texttt{t}, and \texttt{r} is a hypertable on \texttt{t} because that is what buys the compression that Chapter 12 shows has never actually run.

Now recall what \texttt{r.i} is for. Amendment 1 of 2026-07-03 says arrows are reifiable: a relation can be pointed at, and that is what makes the outcome mechanism work. \texttt{r\_active} implements it, and the current definition on the live instance is:

\begin{lstlisting}
CREATE OR REPLACE VIEW r_active AS
SELECT r.* FROM r
WHERE r.d->>'type' <> 'outcome'
  AND COALESCE((SELECT rx.d->>'verdict' FROM r rx
                WHERE rx.d->>'type' = 'outcome' AND (rx.d->>'of')::UUID = r.i
                ORDER BY rx.t DESC, rx.i DESC LIMIT 1), '') <> 'retracted';
\end{lstlisting}

So the retraction mechanism still rests on an unenforced reference into a non-unique column, resolved by a JSON extraction and a runtime cast. What changed is the type and one of the soft links. The cast is \verb|::UUID| now, not \verb|::BIGINT|, and \texttt{retract()} builds \texttt{d->'of'} by \texttt{jsonb\_build\_object} from a \texttt{UUID}, which lands as a JSON string rather than a number. Verified on a throwaway container: insert a witness, retract it, and \texttt{jsonb\_typeof(d->'of')} is \texttt{string}, the cast resolves, and \texttt{r\_active} drops from one row to zero. There is still no foreign key, so an outcome can point at an \texttt{i} that does not exist and the view will simply match nothing. And if two \texttt{r} rows ever share an \texttt{i}, one outcome row retracts both, silently, and \texttt{r\_active} reports exactly the state the retracting author intended for one of them.

The honest reading of what \texttt{uuidv7()} bought here: collisions in \texttt{r.i} are now vanishingly improbable rather than guaranteed on merge, which is a real improvement and is not the same thing as being prevented. The database still cannot tell you that they did not happen.

\begin{Open}{Can \texttt{r.i} be given an enforced uniqueness?}{Settled by one of three: a uniqueness mechanism compatible with hypertable partitioning, dropping the hypertable, or accepting the gap and auditing for it.}
The first route is to exhibit a uniqueness mechanism compatible with hypertable partitioning. A unique index on \texttt{(i, t)} does not suffice, since it permits duplicate \texttt{i} at different \texttt{t}, which is precisely the merge case. The second is to drop the hypertable and measure what the compression policy was actually buying, which per Chapter 12 is currently nothing, since zero chunks have ever compressed. The third is to accept that retraction targets are unenforced and add a periodic audit that counts \texttt{i} values with multiplicity greater than one and reports them as a defect rather than discovering them as wrong answers. The third is cheap and nobody has done it, and it is now cheaper than it has ever been, because the instance it would run against is empty.
\begin{Check}
-- The audit nobody has written. Trivially empty today; not empty after a bad merge.
SELECT i, count(*) FROM r GROUP BY i HAVING count(*) > 1;
-- Why the enforced version is still not available, on a stem-cell container:
ALTER TABLE r ADD CONSTRAINT r_pkey PRIMARY KEY (i);
-- ERROR: cannot create a unique index without the column "t" (used in partitioning)
-- HINT: ... ensure the partitioning column is part of the primary or composite key.
\end{Check}
\end{Open}

\section{What the fix cost, and why the estimate was still wrong}

The whitepaper, section 7, stated the problem correctly and then costed it. Verbatim:

\begin{quote}
One change is required to make this work. Today the internal names for records are counters local to each machine, so two machines would both mint a record numbered five and collide. The fix is to mint names that are unique without any coordination: the fingerprint for anything that is a file or is otherwise a function of its content, and a random, time-ordered identifier for everything else. This is a small, one-time change, and it is the only thing standing between the current single-machine design and painless local copies.
\end{quote}

The diagnosis was right and the prescription was right, and the prescription is what shipped, almost word for word. The adjectives were still wrong, and the way they were wrong is more interesting now that the change has been made than it was when it had not.

\textbf{It was not small, and it was never performed.} The migration the old estimate had to cover was: \texttt{o.i} referenced by \texttt{r.o1} and \texttt{r.o2}, by \texttt{component()}, by \texttt{coauthor}, by \texttt{co\_attended}, by \texttt{r\_active}'s cast; and \texttt{d->'of'} stored as a JSON \emph{number} in 57,994 outcome rows, which a UUID scheme makes type-wrong, so the change was a rewrite of 57,994 JSONB documents rather than an \texttt{ALTER COLUMN}. None of that was done. It was discarded. The Names universe ended on 2026-07-17 and the new one booted with the fix already in it, so the cost of the migration is not a number this project will ever have, because this project declined to pay it and paid something else instead.

\textbf{The primitive did not exist and still does not.} Postgres 17 has no \texttt{uuidv7()}: rechecked on the live container on 2026-07-17, \texttt{PostgreSQL 17.10}, and the function landed in Postgres 18. The whitepaper did not mention this, which is the tell that the change was never costed: you cannot cost a change whose primitive you have not checked exists. It is now costed, at fifteen lines of SQL, sitting in \texttt{stem-cell.sql} above the \texttt{o} table, and the decision it forced is a real one and is now made: the minter lives in the database, not in a client, so anyone with a connection can insert correctly and nobody has to be trusted to bring their own id generator.

\textbf{It was not the only thing.} The one that survives is ordering. The whitepaper's section 10 says ordering across machines must not lean on clock agreement, and \texttt{r\_active} still orders by \texttt{rx.t DESC, rx.i DESC}, where \texttt{t} is \texttt{TIMESTAMPTZ DEFAULT now()}, a wall clock, and \texttt{i} is now a uuidv7, which embeds a wall-clock millisecond. UUIDv7 does not rescue this. It makes the tiebreak deterministic given the timestamps you have, and deterministic is not correct: two machines will agree on the answer and the answer may be the older correction.

This is worth flagging precisely, because \texttt{stem-cell.sql} overstates it. Its identity comment says v7 ordering ``is a free tiebreak that does NOT require trusting that two machines' clocks agree, because nothing downstream reads ordering as truth.'' The second clause is false as written. \texttt{r\_active} reads ordering as truth: latest verdict wins, and latest is decided by \texttt{t} and then by \texttt{i}. The defensible version of the claim is narrower and worth substituting: the ids being time-ordered adds no new dependence on clock agreement, because \texttt{r\_active} already had one through \texttt{t}. The dependence is real, it is old, and it is not this fix's fault, which is a different thing from its not existing.

\begin{Open}{Ordering across machines without trusting clocks}{Settled by two independently minting copies inserting concurrent, conflicting outcomes and \texttt{r\_active} returning the intended verdict, not merely the same one on both.}
The identity blocker is gone, so this is now the blocker. The candidates are the ordinary ones: a logical clock or version vector carried in \texttt{d}, a happens-before edge made explicit as a witness so that ordering is data rather than metadata, or an admission that concurrent retractions of the same witness are a branch under divergence-is-data-v1 and must be realized deliberately rather than resolved by a \texttt{LIMIT 1}. The last is the most consistent with the rest of the design and the least implemented. Nothing short of the two-machine demonstration settles it, because the failure is invisible on one machine, where every clock agrees with itself.
\begin{Check}
-- On the live instance, why this is not settled by uuidv7:
SELECT pg_get_viewdef('r_active', true);   -- ... ORDER BY rx.t DESC, rx.i DESC LIMIT 1
-- t is a wall clock, and the v7 prefix is the same wall clock, to the millisecond:
SELECT version();                                       -- PostgreSQL 17.10
SELECT count(*) FROM pg_proc WHERE proname = 'uuidv7';  -- 1, ours; the built-in is PG18
\end{Check}
\end{Open}

\section{The plan is startable now}

The whitepaper's staged plan was built around a promise, stated in its framing sentence: ``The steps are ordered so that something useful exists before any meeting is needed.'' Step 1, verbatim:

\begin{quote}
The first step needs one person and no meeting. Stand up a lab copy of yoneda from the existing two-table design, which is a clone with nothing to invent, and make the one-time change to globally-unique names so the copy is ready to merge from the start.
\end{quote}

The old objection was that ``a clone with nothing to invent'' and ``make the one-time change to globally-unique names'' sat in the same sentence, and the second half was the thing the document elsewhere admitted was missing and had never been costed, so the plan's zero-cost first step contained its highest-uncertainty item. That objection is spent. The change is made, in the substrate rather than in a copy, and standing up a lab instance is now what the sentence claimed it was: run two files against a fresh container and it is born ready to merge, minus a merge channel to merge over.

The route by which the objection was answered is not the route the plan imagined, and the difference is the whole content of Chapter 4's revision. The plan assumed the identifiers would be rewritten. They were not. The universe holding them was ended, and a new one was booted with the right ones. That is available exactly once, at birth, when there is nothing to lose, and it is not a migration technique.

\section{What this chapter is not saying}

It is not saying the design was wrong. The identity mechanism was the one part of the substrate that was never derived from anything: it is what Postgres hands you when you type \texttt{PRIMARY KEY}, it came in with the first draft in May 2025, and it survived because the system had one writer and one writer cannot see the difference. That is an ordinary way for software to be wrong, and the interesting thing was never the defect but the fourteen months, which are evidence about what this project's methods do and do not catch.

Names are the layer the mathematics does not reach. Everything else in this book was corrected by asking what the structure required. This was not corrected that way, because the structure does not care, and the only thing that would have caught it is a second machine. The second machine finally existed, twice, for about ninety seconds, on 2026-07-17, and it caught it.

And it is not saying identity is finished. \texttt{o.i} is a rigid designator now, which is what \texttt{DESIGN.md} claimed in July 2026 and did not have. \texttt{r.i} is a globally unique name in a column where uniqueness cannot be enforced. The tiebreak that decides which retraction wins is still a wall clock wearing a UUID. One of the three claims the design makes about names is now true, which is one more than yesterday, and it is not three.
